\paragraph{First integrals with isolated critical points.}
The multiplicity argument extends to every fixed value degree.  The
hypothesis below concerns the original critical ideal, including fibers
above repeated roots of the denominator.

\begin{theorem}\label{thm:general-first-integral-zero-dimensional}
Let $b\ge2$, $1\le D<n$, and let $\mathbb K$ have characteristic zero
or characteristic $p>b$.  Suppose $0\ne H\in\mathbb K[X]$ and
\[
 F(X,u)=u^b+\sum_{j=0}^{b-1}a_j(X)u^j,
 \qquad \deg a_j\le(b-j)D.
\]
Assume that the original ideal
\[
 I=(F_u,HF_X-H'F)
\]
is zero-dimensional after extending constants to an algebraic closure;
the unit ideal is allowed.  On any $n$ distinct evaluation coordinates,
every received word has at most
\[
 \left\lfloor b^2+
 \frac{b+3b^2(b-1)^2D/n}{\eta}\right\rfloor
\]
polynomials $P$ of degree at most $D$ satisfying $F(X,P)=cH$ for some
constant $c$ and agreeing with the word on at least $D+\eta n$
coordinates, for every $\eta>0$.  Coordinates and received symbols may
lie in an extension field.
\end{theorem}

\begin{proof}
Extend constants to an algebraically closed field $k$ containing the
coordinates and symbols.  Each section has a unique label $c$, and each
label has at most $b$ sections.  If no nonzero label occurs the claim is
immediate.  Otherwise $N=\deg H\le bD$ by the weighted coefficient caps.
Put $B=HF_X-H'F$ and $\ell=\dim_k k[X,u]/I$.
The bidegrees of $F_u,B$ are at most $((b-1)D,b-1)$ and
$(N+bD-1,b)$.  Homogenize each nonzero generator in
$\mathbb P^1\times\mathbb P^1$ to its own actual bidegree.  Neither
contains a boundary divisor; a common projective curve would therefore
meet the affine chart and contradict zero-dimensionality.  The proper
intersection product gives
\begin{equation}\label{eq:general-critical-length}
 \ell\le(b-1)(N+2bD-1).
\end{equation}
A nonzero constant generator gives length zero directly.

For a section $P$ write $A_P=F_u(X,P)$.  Differentiation gives
$B(X,P)=-HA_PP'$.  Consequently $A_P\ne0$, since otherwise the graph
of $P$ would be contained in the critical scheme, and
\[
 k[X,u]/(I,u-P)\simeq k[X]/(A_P).
\]
After localization, $\operatorname{ord}_x A_P$ is therefore at most the
critical scheme's local length at $(x,P(x))$.

We use one valuation estimate at finite points and infinity.  In a fixed
valued algebraic closure factor
\[
 F_u(X,P+Z)=b\prod_{i=1}^{b-1}(Z-\delta_i),
 \qquad A_P=b\prod_i(-\delta_i),
\]
and choose $\delta$ of maximal valuation.  All $\delta_i$ are nonzero.
If $d=v(A_P)$ then $v(\delta)\ge d/(b-1)$.  Writing the derivative
as $\sum_j c_jZ^j$, each summand of $c_j\delta^j$ is, up to sign,
$A_P$ times a product of ratios $\delta/\delta_i$.  Its valuation is
at least $d$.  Integrating this polynomial in $Z$ thus gives
\begin{equation}\label{eq:general-critical-correction}
 v\bigl(F(X,P+\delta)-F(X,P)\bigr)
 \ge d+v(\delta)\ge\frac{bd}{b-1}.
\end{equation}
Only $1,\ldots,b$ are inverted, so the characteristic assumption suffices.
The argument permits repeated critical roots and ramified valuations.

At infinity this chooses a root of minimal degree.  If
$\deg A_P<(b-1)N/b$, the correction has degree less than $N$ and the
fixed critical value $F(X,P+\delta)/H$ has residue $c$ at infinity.
There are at most $b-1$ fixed critical roots of $F_u$, hence at most
$b-1$ such exceptional labels.  For every other label,
\[
 \deg A_P\ge\lceil(b-1)N/b\rceil.
\]
At a root $x$ of $H$ of multiplicity $m_x$, the inequality
$\operatorname{ord}_x A_P>\lfloor(b-1)m_x/b\rfloor$ similarly forces
one of those critical values to have residue $c$ at $x$.  For each fixed
$x$, at most $b-1$ labels can have this excess.  Fixing one valuation
extension at each place makes this count independent of how the sections
choose critical roots.

Let $r_H$ be the number of distinct roots of $H$, and suppose $r_H>D$.
Select one section for each nonzero label not exceptional at infinity,
and charge it
\[
 C(P)=\sum_{H(x)\ne0}\operatorname{ord}_x A_P+
 \sum_{H(x)=0}\max\left\{0,\operatorname{ord}_x A_P
                    -\left\lfloor\frac{(b-1)m_x}{b}\right\rfloor\right\}.
\]
Off $H=0$ a critical point determines the label $F/H$, so different
selected labels have disjoint charged support.  At a root of $H$ at
most $b-1$ labels have positive excess.  The local length bound gives
$\sum_P C(P)\le(b-1)\ell$.  Also
\begin{align*}
 C(P)&\ge\left\lceil\frac{(b-1)N}{b}\right\rceil
       -\sum_{H(x)=0}\left\lfloor\frac{(b-1)m_x}{b}\right\rfloor\\
 &\ge r_H-\lfloor N/b\rfloor\ge r_H-D>0,
\end{align*}
since each summand in the last sum is at most $m_x-1$.
Adding the exceptional labels and the zero label, and allowing at most
$b$ sections per label, bounds the entire section bank by
\begin{equation}\label{eq:general-root-bank-bound}
 |\mathcal B|\le b^2+
 \frac{b(b-1)\ell}{r_H-\lfloor N/b\rfloor}
 \le b^2+\frac{3b^2(b-1)^2D}{r_H-D}.
\end{equation}
Here \eqref{eq:general-critical-length} and $N\le bD$ give the last step.

Now select any finite list of $L$ sections with at least
$a\ge D+\eta n$ agreements each.  At a coordinate outside $H=0$,
the received symbol fixes one label and matches at most $b$ sections.
Thus
\[
 La\le Lr_H+bn,\qquad r_H-D\ge\eta n-bn/L.
\]
For $L\le\max\{b^2,b/\eta\}$ the stated bound is immediate.
Otherwise $r_H>D$, and \eqref{eq:general-root-bank-bound} implies
\[
 (L-b^2)(\eta-b/L)\le3b^2(b-1)^2D/n.
\]
Expanding and dropping the positive term $b^3/L$ gives the claimed
bound after taking its floor.  Applying this to every finite subset
also excludes an infinite list.
\end{proof}

\begin{corollary}\label{cor:general-first-integral-repeated-fiber}
Keep the weighted coefficient and section hypotheses of
Theorem~\ref{thm:general-first-integral-zero-dimensional}, but do not
assume its critical ideal is zero-dimensional.  Assume characteristic
zero or $p>b(b-1)D$.  If a list exceeds the bound in that theorem, then,
after extending constants to an algebraic closure $k$, there is a
$c\in k$ such that $F(X,u)-cH(X)$ has a repeated nonconstant factor
in $k(X)[u]$.
\end{corollary}

\begin{proof}
The theorem forces a positive-dimensional critical component.  It cannot
be vertical because $F_u$ has nonzero constant leading coefficient $b$.
Its generic root $r$ generates an extension $L/k(X)$ of degree
$d\le b-1$, separable in positive characteristic since $d<p$.
The derivation extends to $L$, and the critical relations give
\[
 \left(\frac{F(X,r)}{H}\right)'=0.
\]
The excessive list includes a nonzero-label section, so $N\le bD$.
The monic weighted equation $F_u/b=0$ makes $r$ integral at finite
places and bounds its pole at each place above infinity by $D$ times
the ramification index.  Therefore $v=F(X,r)/H$ has pole-divisor height
\[
 h_L(v)\le dN+d\max\{bD-N,0\}=dbD\le b(b-1)D.
\]
In characteristic zero a function with zero derivative is constant.
In characteristic $p$, the kernel of this nonzero derivation is $L^p$;
a nonconstant $p$th power has height at least $p$, contradicting the
strict bound.  Hence $v=c\in k$.  The minimal polynomial of $r$ divides
both $F_u$ and $F-cH$, proving the repeated-factor assertion.
\end{proof}

The corollary is a structural condition, not a classification of the
exceptional pencils.  It does not assert that the repeated factor is
linear or that its label descends to the original constant field.
